% Portada adaptada APA 7 y Resumen (páginas 1 y 2)
% evidencia:E03,E08,E13,E18

\thispagestyle{fancy}
\begingroup
\centering
\vspace*{2.5\baselineskip}

{\bfseries MediFlow}\par
\vspace{1.5\baselineskip}

Juan David Campo Mercado\par
Daniel Alejandro Rincón Navarro\par
Carlos David del Castillo Mercado\par
\vspace{0.5\baselineskip}
Universidad del Norte\par
19 de noviembre de 2026\par
\endgroup

\vfill
\begin{flushleft}
{\singlespacing\footnotesize\RaggedRight
\textit{Nota documental.} El título exacto, el orden de autoría y la omisión de campos de asignatura, docente y tutor corresponden a decisión editorial expresa para este informe, adaptando la portada estudiantil de la séptima edición de las normas APA. La fecha indicada es documental y no constituye una fecha oficial de entrega del certamen. Base técnica verificada en el repositorio: commit \nolinkurl{803c1a1f866a978cce8b4622269ccfb4749ba59e}, artefacto web empaquetado \texttt{v2.0} (corte exclusivo \texttt{2026-07-17}, huella SHA-256 \nolinkurl{be8252b056377f8820808a518095235c1215b6ad5c3466a781bad6a0f0ad1e28}).\par}
\end{flushleft}
\clearpage

% evidencia:E08,E13,E18
\section{Resumen}

\noindent MediFlow es una herramienta analítica de apoyo a la planificación operacional del Hospital Universidad del Norte (HUN) construida sobre agregados compatibles con el modelo OMOP CDM v5.4 bajo tres ejes: comprender el flujo asistencial, predecir la demanda diaria y explorar escenarios paramétricos de decisión. Sobre \num{1706001} registros históricos de visitas entre el 1 de enero de 2023 y el último día completo de demanda (\texttt{2026-07-16}; corte exclusivo \texttt{2026-07-17}), la deduplicación SQL consolida \num{1477742} llegadas únicas, el censo por etapas de \texttt{visit\_detail} aplica recorte estadístico $p_{99{,}9}$ y la auditoría de calendario verifica desfase temporal cero. Se documentan por separado dos familias predictivas a 1--7 días evaluadas en 26 pliegues \textit{walk-forward}: el experimento puntual LightGBM Poisson con nivel local y factores externos aceptados (\texttt{lgbm\_final}) reduce el error WAPE diario global de 18,0\,\% en la línea base a 15,6\,\% en regresión de Poisson estratificada y a 10,5\,\% final (mejora relativa archivada de 41,6\,\%; 8,4\,\% en Urgencias y 11,0\,\% en Rehabilitación frente a 43,5\,\% en Tratamiento Urgencias), mientras que el artefacto cuantílico servido en la aplicación web (\texttt{v2.0}, $p_{10}/p_{50}/p_{90}$) registra sobre 182 días un WAPE archivado de 8,91\,\% en Urgencias (cobertura 78,57\,\%), 11,08\,\% en Rehabilitación (64,29\,\%) y 42,18\,\% en Tratamiento Urgencias (73,08\,\%). A 120 días se distinguen los resultados archivados del 3 de noviembre de 2026 (WAPE 15,7\,\%, cobertura 86,5\,\%) de la corrida empaquetada en el artefacto servido. La capa operacional integra dimensionamiento Erlang~C bajo tasas y duraciones supuestas (15 frente a 12 equipos equivalentes en Urgencias 06--12\,h en demanda $p_{90}$; 70 frente a 42 en Rehabilitación), contraste de Ley de Little (\num{277,52} pacientes en población alineada frente a \num{277,47} de media horaria y \num{284,01} en la referencia híbrida del bundle) y simulación horaria de altas tempranas (35\,\% adelantadas reducen las horas rojas de 9,32\,\% a 6,26\,\% manteniendo el pico en 359). Los resultados justifican validación operativa supervisada, pero no demuestran beneficio clínico ni sustituyen la verificación presencial de dotación y camas habilitadas.

\indent\textit{Palabras clave:} flujo hospitalario, OMOP CDM, pronóstico de demanda, validación temporal, Erlang~C, Ley de Little, planificación operacional.

\clearpage
